\section{Buildability includes the means of paying for the build}
\label{sec:economy}

\subsection{The arrangement can be the missing instrument}

A technically feasible construction may remain unused because its beneficiaries are fragmented, its risks are misallocated, or its funding arrives too late. A shared testing facility, a staged acceptance contract, compatible access rights, or a division of responsibilities can be part of the complete means. Such arrangements are construction hypotheses, not inferred agreements.

For participant $i$, a proposed arrangement should compare its expected utility $U_i$ with an explicit outside option $U_i^{\mathrm{out}}$. Modeled voluntary participation requires
\begin{equation}
U_i\ge U_i^{\mathrm{out}}\qquad\text{for each required participant }i.
\label{eq:participation}
\end{equation}
This condition is neither observed consent nor a proof of truthful behavior or coalition stability. It depends on supplied beliefs, information, alternatives, and enforceability. Real acceptance must be obtained at the appropriate stage. New technology can improve competitors' outside options as well as the proposed builder's costs.

\emph{Invent the Economy} makes this joint search explicit \citep{boucherEconomy}. A component with poor stand-alone returns can belong to a worthwhile package. Conversely, positive combined surplus may fail to support acceptable individual terms. The construction brief should identify who benefits, who contributes, when each acts, what each can refuse, and which complete alternative each could choose.

\subsection{Every cash prefix must be feasible}

For a declared cash schedule, write cumulative paid construction and operating costs as $C(t)$, cleared receipts as $R(t)$, received external financing as $F(t)$, financing repayments and distributions as $D(t)$, and required protected funds as $Q(t)$. With opening cash $b$, available cash is
\begin{equation}
A(t)=b+R(t)+F(t)-C(t)-D(t)-Q(t).
\label{eq:cash}
\end{equation}
The route requires $A(t)\ge0$ at every event time. $Q(t)$ must exclude amounts already counted as paid costs or repayments; releases from protection reduce $Q$ rather than create revenue. The equation is a cash-feasibility statement, not a replacement for an operating ledger or an accounting standard.

For a fixed schedule without new financing, the minimum opening cash implied by these events is
\begin{equation}
b_{\min}=\max\left\{0,\sup_{t\le T}\big[C(t)+D(t)+Q(t)-R(t)\big]\right\}.
\label{eq:prefunding}
\end{equation}
A positive terminal contribution does not remove an earlier funding shortfall. A promised investment is not received financing. Refund exposure, supplier commitments, and support obligations can remain after collection. A lower complete cost is useful only when the revised sequence can actually occur.

Starting capital is a design variable for a new economic comparison. Small staged construction, outside financing, a larger initial facility, and non-entry may all be legitimate alternatives. Their terms, costs, downside, and timing belong in the comparison. An earlier fixed-capital experiment retains its original assumptions; it is not silently topped up to make a new result attractive.

\subsection{What the finite study says about continuation}

The study uses synthetic planning units, with no currency interpretation. Its source route costs 54 and produces an assumed receipt of 80 after accepted delivery. Preserving the fixture alone leaves a successor cost of 44. Acquiring the supplied reusable constructor costs 9; each successor then costs $33+m$, where $m$ is its current maintenance obligation. Thus, for $n$ successors,
\begin{align}
C_{\mathrm{fixture}}(n)&=54+44n,\\
C_{\mathrm{constructor}}(n,m)&=54+9+n(33+m),\\
\Delta C(n,m)&=n(11-m)-9.
\label{eq:retention}
\end{align}
Constructor retention is strictly cheaper exactly when $n(11-m)>9$, conditional on the same acceptance, valid current rights, deadlines, fixture availability, and feasible cash prefixes. At $m=0$, three successors save 24 units after acquisition. At $m=12$, every additional successor makes this cost comparison worse. The parameters illustrate a threshold; they are not observed market costs or evidence of increasing returns in a real enterprise.

The distinction between \emph{later}, \emph{financed by}, and \emph{enabled by} is visible even here. An illustrative cash path starting at 54 spends the source cost, collects the modeled 80, pays 9 for acquisition, and can prepay the next 33-unit delivery. Yet the source's net contribution is only 26. If all opening capital must remain protected, the 26 does not cover acquisition plus the next 33. The path is feasible using available modeled cash; it is not, on that rule, fully funded by unrestricted earnings alone.

Likewise, a retained constructor can make a 33-unit successor affordable where a 44-unit fixture-only route is not. To establish this as a \emph{causal resource advantage}, compare equally resourced successors that may build or buy alternative means. To establish the value of acquiring the constructor, charge its acquisition and prior failed development as well. State-level reachability and whole-lineage investment value answer different questions.

\subsection{Value creation and value capture must remain connected}

An enterprise can use accepted value to sustain future construction through prices, subscriptions, licensing, shared facilities, or other voluntary arrangements. The form should reward the useful outcome and bear the necessary obligations. Billing unnecessary operations can conflict with an intelligence whose best contribution is to eliminate those operations.

Complete economic cost includes search, unsuccessful candidates, development, acquisition, production, checking, integration, maintenance, support, financing, and residual liabilities. The costing policy should explain amortization and allocation without counting one expense twice or crediting the same receipt to several generations. Included assistant resources do not imply zero underlying economic cost; unavailable bills remain unknown.

Commercial recurrence need not prove scientific novelty. A known method can support an excellent business. Scientific construction need not immediately produce a sale. The desired flywheel joins useful capability, voluntary accepted delivery, resources available for continuation, and a tested successor contribution; each connection retains its own evidence.
